% Source: MQ33a_v1.html
\documentclass[11pt]{article}
\usepackage[utf8]{inputenc}
\usepackage[T1]{fontenc}
\usepackage[margin=1in]{geometry}
\usepackage{amsmath}
\usepackage{amssymb}
\usepackage{siunitx}
\usepackage{enumitem}
\usepackage{parskip}

\title{PS160 -- MQ33a\_v1.html}
\date{}

\begin{document}
\maketitle

\section*{Questions}

\begin{enumerate}[leftmargin=*]
  \item Suppose two mirrors are arranged perpendicular to each other, and a ray of light is incident on the vertical mirror at an angle of 58$^{\circ}$ with respect to the normal, as in the figure. Find the angle $\phi$ with respect to the normal.

[Image: ME33\_plane\_mirror\_3.png]

$\phi$ = $^{\circ}$ \hfill \textit{(1 pts, Numeric)}

  \item Light with a wavelength of 536 nm in vacuum enters a material with an index of refraction of 2.8. Calculate its wavelength in nm in the medium.

$\lambda$ = nm \hfill \textit{(1 pts, Numeric)}

  \item Light with a vacuum wavelength of 502 nm travels at 2.01 $\times$ 10\textsuperscript{8} m/s in a certain medium. Calculate its wavelength in nm in the medium.

$\lambda$ = nm \hfill \textit{(1 pts, Numeric)}

  \item A ray of light traveling in a fluid with index of refraction 1.4 is incident at an angle of 39$^{\circ}$\textsuperscript{}with respect to the normal on a transparent material with index of refraction 2.3. Calculate the refraction angle in degrees. \hfill \textit{(1 pts, Numeric)}

  \item The figure shows a light ray incident on a pane of glass in air with $\theta$ = 33$^{\circ}$. Calculate the lateral distance \emph{d} in cm the ray is shifted if \emph{n}\textsubscript{glass} = 7 and \emph{h} = 8 cm.

[Image: ray\_shift\_01.png]

\emph{d} = cm \hfill \textit{(1 pts, Numeric)}

\end{enumerate}

\end{document}